% Place after the bibliography.
% Required packages: booktabs, array.
% Replace the previous appendix block; do not append this block to it.
% Use \appendix only once in each language's manuscript.

\clearpage
\appendix

% Appendix A
\section{Reproducibility and Measurement}
\label{app:protocol}

\paragraph{Reproduction materials.} The GitHub release described in the
reproducibility statement will provide the executable source, configuration
and dependency manifests, run commands, figure-generation scripts, and
machine-readable result tables needed to regenerate the paper. Missing or
unevaluable fields will be reported explicitly rather than imputed.

\paragraph{Dataset provenance and naming.} We use the public
\emph{HumDial-FDBench} test archive, whose dataset paper is Wang et al.
\citep{humdial2026} and whose data page is cited separately
\citep{humdialdataset2026}. The official split is 8,898 Train, 1,800 Dev,
and 5,000 Test instances across nine scenarios; the public test archive also
contains 4,100 clean scoring files. Its pinned revision is
\texttt{ed1770a6265ffe9f0e5cd999290936ccab43482b}, and the Hugging Face data
card labels it Apache-2.0. In the pinned archive, WAV payloads are mono,
16-bit PCM at 16\,kHz. We use the released WAV bytes and JSON annotations as
supplied, with no project resampling, denoising, loudness normalization, or
synthetic mixing. Source IDs retain the archive path
\texttt{test/\{en,cn\}\_test\_nondev/\{scenario\}/\{base\_id\}[\_add].wav};
clean files are separate scoring references. A runtime cut is generated from
the complete selected WAV at the fixed nonempty state boundary used by the
experiment, so it is a project-derived replay cut rather than an official
HumDial-FDBench segment or leaderboard score; cut IDs, timestamps, members,
and hashes are bound in the run manifest. The release is human-recorded
performed speech, not TTS audio. Because the public paper and data card do
not document de-identification or speaker anonymization, we make no claim
that the audio is de-identified.

Throughout, \emph{FD-Bench} (Peng et al.) denotes a separate benchmark and
is never shortened to the unhyphenated name; ``HumDial-FDBench'' is retained
only as the official HumDial dataset identifier. \emph{Full-Duplex-Bench} denotes Lin et
al.'s family: v1 and v1.5 are offline releases, while
\emph{Full-Duplex-Bench-v2} is the later real-time framework. \emph{HumDial-FDBench}
is the HumDial benchmark built on the v1.5 protocol; it is not the FD-Bench
dataset and not Full-Duplex-Bench-v2. Our adapted event definitions are
described as ``FD-Bench-style''.


The cold-resumption study uses the same Lychee-FD model and local
Token2Wav backend \citep{lychee2026} on two A100-SXM4-40GB GPUs,
with model execution and one acoustic lane placed separately.
Shared weights remain resident.
Model, sampling, numerical, input, and playback settings are matched
within each comparison. Reproduction requires the executed source,
dependency, and configuration manifests, including archived patches,
rather than a repository commit alone.

The final matrix contains 96 executions from 12 source groups,
four policies, and two repetitions. E1 contains 44 executions and
E2 contains 52; their measurements are analyzed separately.
The measured cohort includes 24 pending speech tokens and
20,032 buffered PCM samples at 24\,kHz. These values describe the
predefined nonempty cut, not a natural conversation-state distribution.
The nominal suspension is 30 seconds, but latency uses actual events.
New-audio rendering spans Worklet resumption to the first-new-PCM frame
rendered within the same AudioContext; buffered
playback is not new computation. The primary metric is
\[
T_{\mathrm{first\text{-}new\text{-}PCM}}
=t_{\mathrm{first\ new\ PCM\ rendered}}
-t_{\mathrm{Worklet\ resumed}}.
\]
Save is the interval from pause acceptance to an independent checkpoint
becoming ready. Ready is the interval from RESUME acceptance to execution
becoming ready. Gap is the interval from buffered PCM exhaustion to the
first-new-PCM frame rendered by the AudioContext.
These endpoints are not serial segments and must not be added.

Configuration summaries aggregate repetitions with the corresponding
endpoint within each source and environment before taking the median
across sources. Paired differences first match source, repetition, and
environment; missing values are not replaced with zero. The cold,
transfer, throughput, and interaction cohorts retain separate identities
and are not pooled.

The final matrix contains 92 technical passes out of 96 scheduled runs.
Rebuild has one retained E1 cancellation after native-browser frame
overlap; Hybrid has three retained E1 failures (two frame-overlap
cancellations and one frozen-cut/first-PCM observation-round mismatch).
Hot and Joint have no invalid runs, and all E2 runs pass the measured
technical contracts. Hot and Rebuild have no corresponding checkpoint,
pinned-buffer, or GPU-workspace measurement events; these cells are marked
N/A or NM in the main table rather than filled with zero. System pinned,
$L_{\mathrm{useful}}$, and $J(T)$ are outside the measured cold-resumption
contract and remain explicitly unevaluated.


% Appendix B
\section{Joint State and Resumption Validation}
\label{app:state_contract}

Table~\ref{tab:app_state_contract} summarizes the joint resumption
boundary. Computational state is restored, whereas the external ledger
of accepted input, cancellation decisions, and delivery records remains current.
This distinction applies consistent-snapshot and rollback-recovery
principles at request granularity \citep{chandy1985,elnozahy2002}.

\begin{table}[tbp]
    \centering
    \small
    \setlength{\tabcolsep}{5pt}
    \caption{
        State coverage and resumption obligations.
        External records are retained rather than overwritten
        with an older checkpoint copy.
    }
    \label{tab:app_state_contract}
    \begin{tabular}{
        @{}
        >{\raggedright\arraybackslash}p{0.20\linewidth}
        >{\raggedright\arraybackslash}p{0.73\linewidth}
        @{}
    }
        \toprule
        State class & Contents and resumption obligation \\
        \midrule
        Model
        & Effective KV, non-KV histories, processor/control state,
          sampling context, and logical positions.
          Restore values in legal mappings and rebuild active bindings. \\
        Bridge
        & Produced but unconsumed tokens, partial buffers,
          initialization/flush state, and pending output.
          Preserve content and consumption order. \\
        Acoustic
        & StreamingDecoder and Token2Wav continuation,
          Flow / HiFT caches, actual random context, and pending PCM.
          Restore a compatible execution context. \\
        Identity and progress
        & Request, generation, state version, cut, lease epoch,
          sequence, producer epoch, and stage watermarks.
          Validate ownership against the current lifecycle. \\
        External records
        & Accepted input, cancellation, and delivery/rendering intervals.
          Preserve current facts and reject obsolete submissions. \\
        \bottomrule
    \end{tabular}
\end{table}

Capture establishes a consistent cut and an independent host
checkpoint before releasing request-private resources.
Resumption prepares legal mappings, rebuilds bindings, and validates
the participating states before enabling execution or output.
Current generation, version, and cancellation state are rechecked
at publication; a failed or cancelled resumption is not published.
The restored input-consumption position does not overwrite input
received after the checkpoint.

The request ID regression distinguishes physical output order
$[A,B]$ from a sampler selection containing only $B$.
request ID lookup retrieves $B$'s corresponding side outputs rather
than physical row zero.
Separate archived checks cover pending bridge content,
pre-publication rejection, and request-local cancellation.
Complete token and waveform comparisons are conditioned on matched
input/control histories and numerical configurations; they do not
establish semantic correctness or unconditional kernel determinism.

\paragraph{Protocol proof obligations.}
The protocol envelope $q$ contains the five fields
\textit{request\_id}, \textit{generation}, \textit{state\_version},
\textit{cut\_id}, and \textit{lease\_epoch}; each record extends it with
\textit{sequence} and \textit{producer\_epoch}.  The following proof sketch
states the induction used by the implementation.  The base state has one
owner and empty produced/consumed sets.  SAVE closes all writers before the
cut linearization point, so every record in the checkpoint has the same $q$;
draining in-flight commits makes the checkpoint and queue partition the
produced-but-unconsumed set.  RESTORE installs a fresh lease epoch only after
the current ledger watermarks have been read.  Each subsequent commit is
accepted only when its complete envelope equals the installed permit, and
the ledger update and queue-position advance are atomic.

Consequently, ownership is preserved because a record with another request ID
is rejected; bridge conservation is preserved because a token is either in
the checkpoint or removed exactly once by the consumer transaction; and ledger
monotonicity is preserved because input, cancellation, and delivery watermarks
only advance.  The generation, state version, and lease epoch together reject
late workers and ABA lease reuse.  The consumer accepts PCM only at the next
sequence number, which proves the prefix/order invariant by induction.  Since
PUBLISH is enabled only after these checks and is the sole transition out of
RESTORING, no invalid state can execute or become visible.  If any premise is
false, the transition is to \textsc{Aborted}; hence failure cannot publish a
partial restore or roll back an external fact.

The two linearization points are observable in the trace: \textsc{Cut} is the
atomic close-and-snapshot event, and \textsc{Restore} is the atomic install of
the new lease and permit.  Cancellation linearizes at its ledger append; a
resume that observes that append aborts, while a later cancellation invalidates
the permit before the next commit.  These rules cover mismatched cuts, stale
versions, ABA leases, late workers, duplicate bridge consumption, reordered
PCM, cancellation/resume races, and post-checkpoint input without relying on
physical worker identity.

\begin{algorithm}[tbp]
\caption{Save/resume protocol for logical request $r$.}
\label{alg:save_resume}
\begin{algorithmic}[1]
\Require Current envelope $q$, states $(M,B,A)$, external ledger $L$
\Ensure Either a published permit or an aborted, non-executing request
\State Close intake, token production, and output publication atomically
\State Increment \textit{state\_version}; allocate \textit{cut\_id}
\State Quiesce $M,B,A$ and drain in-flight commits at a supported boundary
\If{any record has an envelope different from $q$}
    \State reject the record and abort the capture
\EndIf
\State $K \gets \Call{Snapshot}{M,B,A,L,q}$ at the \textsc{Cut} linearization point
\If{$\neg\Call{CheckJoint}{K}$}
    \State retain the original resources; \Return \textsc{Aborted}
\EndIf
\State Release private resources and mark $r$ suspended
\State Read current $L$; reject if cancelled or below its input/cancel/delivery watermarks
\State Allocate a fresh \textit{lease\_epoch} and merge post-cut input
\State Rebind $M,B,A$ from $K$
\If{$\neg\Call{CheckJoint}{M,B,A,L,q}$}
    \State discard unpublished state; \Return \textsc{Aborted}
\EndIf
\State Install the new lease and permit atomically at \textsc{Restore}
\While{$r$ is executing}
    \If{commit envelope $\ne q$ or PCM sequence is not the next prefix value}
        \State discard the commit
    \Else
        \State atomically append the ledger record and advance the consumer position
    \EndIf
\EndWhile
\State \Return \textsc{Published}
\end{algorithmic}
\end{algorithm}


% Appendix C
\section{Additional Performance Results}
\label{app:performance}

Table~\ref{tab:app_latency} supplements the main timing results.
It reports Save, Ready, the first-new-PCM render interval, and Gap with
their own anchors. Each interval retains its own endpoints and is not
added to another interval to construct a new end-to-end metric.

\begin{table}[tbp]
    \centering
    \small
    \setlength{\tabcolsep}{6pt}
    \caption{
        Additional cold-resumption measurements in seconds.
        Values are source-level medians within each environment.
        Each value is a source-level median for the corresponding endpoint;
        N/A identifies an inapplicable event and NM identifies a missing
        measurement event; neither is zero. The first block reports the
        latest policy-level validity over the full 96-run denominator;
        failures remain counted in that denominator.
    }
    \label{tab:app_latency}
    \begin{tabular}{@{}llcrrrr@{}}
        \toprule
        Environment & Policy & Valid/total & Save & Ready & First-new-PCM render & Gap \\
        \midrule
        \multicolumn{7}{@{}l}{\emph{All environments (latest 96-run cohort)}} \\
        All & Hot     & 96/96 & \na & \na & \na & \na \\
        All & Rebuild & 92/96 & \na & \na & \na & \na \\
        All & Joint   & 96/96 & \na & \na & \na & \na \\
        All & Hybrid  & 96/96 & \na & \na & \na & \na \\
        \midrule
        E1 & Hot     & 11/11 & 0.004 & 0.001  & 1.184 & 0.349 \\
        E1 & Rebuild & 10/11 & 1.847 & 13.317 & 14.451 & 13.616 \\
        E1 & Joint   & 11/11 & 1.387 & 0.699  & 1.899 & 1.064 \\
        E1 & Hybrid  & 8/11  & 0.655 & 0.130  & 1.349 & 0.515 \\
        \midrule
        E2 & Hot     & 13/13 & 0.004 & 0.001  & 1.168 & 0.333 \\
        E2 & Rebuild & 13/13 & 1.926 & 13.250 & 14.400 & 13.565 \\
        E2 & Joint   & 13/13 & 1.387 & 0.700  & 1.872 & 1.037 \\
        E2 & Hybrid  & 13/13 & 0.696 & 0.136  & 1.315 & 0.480 \\
        \bottomrule
    \end{tabular}
\end{table}

In the final cohort, the first-new-PCM render interval exceeds the
post-buffer gap by $20{,}032/24{,}000=0.834667$ seconds. This is the
duration of the predefined buffered playback, not an additional latency
term. The two values are related views of the same delivery timeline,
not independent improvements.
In E2, Joint and Hybrid hold independent checkpoint payloads of
672.207 and 22.416\,MiB, respectively, and managed pinned buffers
of 256 and 128\,MiB.
Both retain model-device workspaces, producing a net allocation
increase of 60.489\,MiB.
Checkpoint size, pinned storage, process RSS, and device allocation
are distinct, non-additive measurements.

The independent transfer study contains 12 Sync/Optimized pairs
from six source groups.
The transfer-to-ready relative reduction is calculated within each pair as
$(T_i^{\mathrm{sync}}-T_i^{\mathrm{opt}})/T_i^{\mathrm{sync}}$.
Its median is 16.00\%, and the median absolute reduction is
0.533 seconds, with all 12 pairs favoring the optimized path. This cohort
does not measure Worklet resumption, first-new-PCM rendering, playback
completion, or end-to-end resumption latency.
The additional paused and sampled peak RSS are median paired
243.2 and 182.2\,MiB, respectively.
The archived throughput matrix instead averages within-repeat
ratios for each workload/concurrency cell.
Its full set of configurations and repetitions must be retained;
new throughput measurements form a separate cohort.


% Appendix D
\section{Interaction and External Protocols}
\label{app:interaction}

The natural-interaction pilot uses adapted FD-Bench-style definitions
\citep{fdbench2025}.
SRR retains eligible non-interruption questions, SIR actual
interruption opportunities, and SRIR successfully interrupted events
as their respective denominators.
EIR records premature assistant entry.
FSED spans user-utterance end to matched reply playback onset;
IRD spans interruption onset to the end of preceding assistant speech.
Signed timings are retained.
Explicit CPU suspension is not counted as a speech interruption,
and PCM concatenation does not replace the actual rendering timeline.

Operationally, the exported event columns compute
$\mathrm{SRR}=k_{\mathrm{SR}}/n_{\mathrm{SR}}$,
$\mathrm{SIR}=k_{\mathrm{SI}}/n_{\mathrm{SI}}$,
$\mathrm{SRIR}=k_{\mathrm{SRI}}/n_{\mathrm{SRI}}$, and
$\mathrm{EIR}=k_{\mathrm{EI}}/n_{\mathrm{EI}}$; a missing column value is
excluded from that metric's denominator. The controlled path has
10 eligible questions, 3 realized interruption opportunities, and
0 invalid/cancelled runs out of 6 per path. The only nonzero failure-like
field is one Omni reply timeout, which removes one FSED observation and
leaves $n_{\mathrm{FSED}}=9$ for that path. Percentile intervals in the
main table are source-level bootstrap intervals (200,000 resamples of the
six source groups, seed 20260926); the valid-bootstrap count is retained in
\texttt{interaction\_metrics\_bootstrap.csv}.

The shared Lychee/DuplexPilot online path is one implementation,
not two independent baselines.
It uses 400\,ms end-of-chunk input delivery, whereas the Omni client
uses 200\,ms send-then-wait delivery.
Models and actual interruption exposure also differ.
The pilot therefore describes application behavior rather than
isolating planned-cold-resumption effects.
Prior listening feedback applies only to the reviewed examples and
is not extended to unreviewed pilot outputs.

Two native vLLM-Omni hot reattachments \citep{omnidoc} produce
request-to-ACK times of 1.512 and 1.421\,ms and first-PCM reception
times of 725.412 and 776.220\,ms.
Recorded PCM delivery has no missing or duplicate sequence events
in these cases, but the production phase of the first returned PCM
is unknown.
The LiveServe-style preloading component \citep{liveserve2026}
prepares 20 and 15\,MiB of KV approximately 9.9 seconds in advance.
It improves readiness by approximately 56.65 and 100.01\,ms,
while new-PCM reception from the common source-window anchor is
later by 23.7 and 36.4\,ms.
These observations retain their native protocols and do not supply
a unified external planned-cold-resumption comparison.

\begin{table*}[t]
\centering
\scriptsize
\setlength{\tabcolsep}{3pt}
\caption{Per-source-repeat interaction results for the controlled Lychee/DuplexPilot and native Omni reference paths. Rates are reported as numerator/denominator; FSED and IRD are signed playback-layer medians in milliseconds with valid $n$. Every repeat completed without an invalid or cancelled run. The backchannel exclusions are intentional non-metric rows, not failures.}
\label{tab:interaction_repeats}
\begin{tabular}{llccccrrcc}
\toprule
Config & source group & SRR & SIR & SRIR & EIR & FSED (ms) & IRD (ms) & timeout & inv/can \\
\midrule
Lychee & 0001\_0001 & 1/1 & -- & -- & 0/1 & 3666.4 (1) & -- & 0 & 0 \\
Lychee & 0001\_0003 & 1/1 & -- & -- & 0/1 & 3762.9 (1) & -- & 0 & 0 \\
Lychee & 0001\_0015 & 1/1 & 1/1 & 1/1 & 1/2 & 5582.8 (2) & 266.6 (1) & 0 & 0 \\
Lychee & 0001\_0025 & 1/1 & 1/1 & 1/1 & 1/2 & 5575.1 (2) & 278.2 (1) & 0 & 0 \\
Lychee & 0001\_0029 & 2/2 & -- & -- & 1/2 & 5808.0 (2) & -- & 0 & 0 \\
Lychee & 0001\_0031 & 1/1 & 1/1 & 1/1 & 1/2 & 3615.5 (2) & 898.9 (1) & 0 & 0 \\
\midrule
Omni & 0001\_0001 & 1/1 & -- & -- & 1/1 & -185.9 (1) & -- & 0 & 0 \\
Omni & 0001\_0003 & 1/1 & -- & -- & 1/1 & 1324.7 (1) & -- & 0 & 0 \\
Omni & 0001\_0015 & 1/2 & -- & -- & 2/2 & 302.9 (1) & -- & 1 & 0 \\
Omni & 0001\_0025 & 1/1 & 1/1 & 1/1 & 1/2 & 1507.4 (2) & 3444.6 (1) & 0 & 0 \\
Omni & 0001\_0029 & 1/1 & 1/1 & 1/1 & 2/2 & 1351.6 (2) & 564.1 (1) & 0 & 0 \\
Omni & 0001\_0031 & 1/1 & 1/1 & 1/1 & 2/2 & 1544.5 (2) & 2068.3 (1) & 0 & 0 \\
\bottomrule
\end{tabular}
\end{table*}

